%!TEX root = main.tex
\chapter{Transformatoren}
\label{Transformatoren}

Het schema voor de transformer kan op de volgende manier worden getekend:

\begin{center}
    \begin{circuitikz}[scale=0.8] 
        \draw (-1,-2) to[isourcesin] (-1,2);
        \draw (-1,2) to[short] (2,2);
        \draw (2,2) to[short] (2,1);
        \draw (3,1) to[short] (1,1);
        \draw (1,1) to[cute inductor, l_=$jX_c$] (1,-1);
        \draw (3,1) to[american resistor, l_=$R_c$, i_=$I_0$] (3,-1);
        \draw (3,-1) to[short] (1,-1);
        \draw (2,-1) to[short] (2,-2);
        \draw (2,-2) to[short] (-1,-2);
        \draw (2,2) to[short] (4,2);
        \draw (4,2) to[L, l=$jX_1$, i_=$I_1$] (6,2);
        \draw (6,2) to[american resistor, l_=$R_1$] (9,2);
        \draw (9,2) to[cute inductor] (9,-2);
        \draw (9,-2) to[short] (2,-2);
        \draw (10,2) to[cute inductor] (10,-2);
        \draw (10,2) to[L, l=$jX_2$, i_=$I_2$] (12,2);
        \draw (12,2) to[american resistor, l_=$R_2$] (14,2);
        \draw (14,2) to[european resistor, l_=$R_L$] (14,-2);
        \draw (14,-2) to[short] (10,-2);
        \draw (8,2) to [open, v^>=$V_1$] (8,-2);
        \draw (11,2) to [open, v^>=$V_2$] (11,-2);
    \end{circuitikz}
\end{center}

Binnen de schakeling zijn 2 verliezen te berekenen. Koper verliezen ($P_{Cu}$) en ijzer verliezen ($P_{Fe}$). 

$P_{Cu} = I_1^2 \cdot R_1 + I_2^2 \cdot R_2$

$P_{Fe} = \frac{V_s^2}{R_c}$

Het ijzer verlies is vanwege de kern. Daarom wordt het ook geschreven als $P_c$ (core).

Naast de koper verliesen en ijzer verliezen, is er nog een klein beetje extra verlies, dit heet hysterese verlies.

Voor het rekenen aan de bovenstaande schakeling zij de volgende formules van belang:

$\frac{V_2}{V_1}=\frac{N_2}{N_1}$

$\frac{I_2}{I_1}=\frac{N_1}{N_2}$

\vspace{5mm}

Ook geldt:

$I_0$= nullasttroom

\vspace{5mm}

In deze formules is $N$ het aantal windingen van de spoel. Of de verhouding van de windingen.

$\varepsilon = N \cdot \frac{d\Phi}{dt}$

\vspace{5mm}

in de schakeling geld: 

$P_1=P_2$

Dus

$V_1 \cdot I_1 = V_2 \cdot I_2$

\vspace{5mm}

De percent regulation kan worden berekend met de volgende formule:

$\frac{V_{no-load}-V_{load}}{V_{load}} \cdot 100\%$

\vspace{5mm}

De power efficiency kan worden berekend met de volgende formule:

$\eta = \frac{P_{load}}{P_{in}}$

$P_{load} = P_{in}-P_{loss}$

$P_{loss} = P_{Cu} + P_{Fe}$

\vspace{5mm}

Om het rekenen aan de schakeling te vereenvoudigen kan je reflecteren naar links. Dan komt de schakeling er als volgende uit te zien:

\begin{center}
    \begin{circuitikz}
        \draw (0,0) to[isourcesin] (0,2)
        to[american resistor, l=$R_1$] (4,2)
        to[european resistor, l=$Z_L'$] (4,0)
        to[short] (0,0);
        \draw (3,2) to[open, v^>=$V_1$] (3,0);
    \end{circuitikz}
\end{center}

Hier geldt:

$Z_L' = \left( \frac{N_1}{N_2} \right)^2 \cdot Z_L$

\section{Nullastproef}
De nullastproef wordt gebruikt om de ijzer verliezen te bepalen. Hier wordt geen load aangesloten. De schakeling ziet er als volgt uit:

\begin{center}
    \begin{circuitikz} 
        \draw (-1,-2) to[isourcesin, l^=$U_{nom}$] (-1,2);
        \draw (-1,2) to[short, -*] (4,2);
        \draw (2,2) to[short] (2,1);
        \draw (3,1) to[short] (1,1);
        \draw (1,1) to[cute inductor, l_=$jX_c$, i_=$I_\mu$] (1,-1);
        \draw (3,1) to[american resistor, l_=$R_c$, i_=$I_{Fe}$] (3,-1);
        \draw (3,-1) to[short] (1,-1);
        \draw (2,-1) to[short] (2,-2);
        \draw (-1,-2) to[short, -*] (4,-2);
    \end{circuitikz}
\end{center}

Hier geldt:
$U_{nom}$ = nominale spanning
$I_\mu$ = magnetiseringsstroom
$I_{Fe}$ = ijzer verlies stroom
$P_{Fe}$ = vermogen
 
$P_{Fe} = \frac{U_{nom}^2}{R_c}$

$R_c=\frac{N_{nom}^2}{R_L}$

$I_{Fe} = I_c=\frac{P_{Fe}}{U_{nom}}$

$I_0 = \sqrt{I_\mu^2 + I_{Fe}^2}$  

$X_c=\frac{U_{nom}}{I_\mu}$

\section{Kortsluit proef}
De kortsluit proef wordt gebruikt om de koper verliezen te bepalen. Hie wordt het kort gesloten. De schakeling ziet er als volgt uit:

\begin{center}
    \begin{circuitikz}
        \draw (0,0) to[isourcesin, i^=$I_{nom}$] (0,2)
        to [L, l=$jX_{k}$] (2,2)
        to [american resistor, l=$R_k$] (4,2)
        to [short] (4,0)
        to [short] (0,0);
    \end{circuitikz}
\end{center}

Hier geldt:

$I_{nom}$ = nominale stroom

$P_{Cu} = I_{nom}^2 \cdot R_k$

$X_k=X_1+X_2'$

$X_k= \sqrt{Z_k^2-R_k^2}$ %let op verbetert naat -

$R_k=R_1+R_2'$